% Résumé du mémoire.
% Abstract in French.
%
\chapter*{ABSTRACT}\thispagestyle{headings}
\addcontentsline{toc}{compteur}{ABSTRACT}

Cars spend 95\% of their time in a parking space and a growing literature suggests that the free provision of parking plays a role in travel behaviour. Municipal governments have reacted to any issue surrounding parking by requiring that spaces be constructed with any new development through minimum parking regulations.\par
This thesis proposes a web site and attending infrastructure to store and disseminate information concerning parking to reduce the fragmentation and lack of coverage of data concerning parking supply.\par
The tool and database are divided into four key sections. The first allows the storage and visualization of different estimates of the parking supply for any parcel on a territory. The second allows the storage and modelling of minimum parking requirements as a function of time, space and land use and to model the amount of parking required for every tax roll entry with valid data. The third creates a framework for selection of a random sample of parcels for quantitive assessment of the difference between built parking and the minimum requirements.  The fourth is an aggregation, visualization and analysis tool of the prediction results.

Quebec City is presented as a case study. The regulations of the city are traced backed to the 1990s and the regulations of that period are applied to any property built before. The geopolitics of the city are modelled back to the 1960s.\par

The model predictions are assessed against a random sample of properties across multiple land uses. Graphical analysis and bootstrap methods are used to calculate error intervals on mean error and mean absolute error for diagnostics. They reveal systematic biases for single family residential properties and sizeable errors for all categories, in part explained by the unit conversion process required to apply certain parking minimums.

Finally, the manual additions made to the parking supply estimates are discussed before proposing different indicators to assess parking supply and demand.

Given the errors shown in the assessment phase, methods for missing data prediction, improvements to unit conversions and the robustness of the regulation data model are suggested before using the model in a decision-making context. Future research should deploy survey based inference methods to build a robust counterfactual to deductive methods commonly used to assess parking supply.